% PILOTO EQUILIBRADO: se conserva el itinerario completo del tema original.
% 21 diapositivas originales (mismos ejemplos, ecuaciones y 13 imágenes);
% 14 diapositivas nuevas de contexto, transición, interpretación y práctica.
% Las diapositivas clásicas se amplían ligeramente cuando faltan explicaciones.
% Las cuatro imágenes finales conservan su protagonismo y disponen de
% una interpretación matemática en una diapositiva inmediatamente posterior.
% Los identificadores "Diapositiva original N" facilitan la comparación.
% Archivo independiente: no modifica tema1.tex, main.tex ni otros pilotos.
% Compilar desde la raíz del repositorio.
\documentclass[main.tex]{subfiles}

\begin{document}
\graphicspath{{imagenes/tema1/}{../imagenes/tema1/}}

% Diapositiva original 1
\portadatema{1}{Funciones de varias variables}
% NUEVA EQUILIBRADA: Índice sin portadas adicionales
\begin{frame}{Contenido del tema}
\begin{enumerate}
\item Funciones de varias variables: definición y ejemplos.
\item Dominio, imagen y representación gráfica.
\item Determinación de dominios.
\item Operaciones con funciones.
\item Curvas de nivel: definición, cálculo e interpretación.
\item Aplicaciones y ejercicios.
\end{enumerate}
\medskip
\small
Partiremos de conceptos conocidos del cálculo de una variable y
los generalizaremos a funciones de varias variables.
\end{frame}

% NUEVA EQUILIBRADA: Contexto matemático para las diapositivas clásicas
\begin{frame}{De una variable a varias variables}
Hasta ahora hemos estudiado funciones que dependen de una sola variable:
\[
f:D\subseteq\mathbb R\longrightarrow\mathbb R,\qquad y=f(x).
\]
\medskip
En numerosos problemas científicos y de ingeniería, la magnitud que
queremos calcular depende de \textbf{varios datos de entrada}.
\medskip
\[
T=T(x,y),\qquad V=V(x,y,z).
\]
\medskip
La idea de función es la misma: a cada entrada válida se le asigna
un \textbf{único valor de salida}. Cambia el número de variables
independientes.
\end{frame}

% Diapositiva original 2
\begin{frame}{Ejemplo: temperatura sobre una superficie}
La temperatura $T$ en un punto de la superficie de la Tierra a una hora determinada depende de la longitud y la latitud de dicho punto.
\medskip

Se puede pensar en $T$ como una función de dos variables y, si las denominamos $x$ e $y$, esta dependencia funcional se puede poner como
\[T=f(x,y).\]
\end{frame}
% Diapositiva original 3
\begin{frame}{Ejemplo: volumen de una caja}
El volumen de una caja rectangular de dimensiones $x,y,z$ es $\text{Volumen}=x\cdot y\cdot z$.
\medskip

Este es un ejemplo de una función real de tres variables reales, que simbolizamos por:
\[f(x,y,z)=x\cdot y\cdot z.\]
\figura{s3_2.png}{.45}{.43}
\end{frame}
% Diapositiva original 4
\begin{frame}{Definición de función de varias variables}
\textbf{Definición}
\medskip

En general, una función real de $n$ variables reales es una correspondencia que asigna a cada $n$-tupla $(x_1,x_2,\ldots,x_n)$ de su dominio un único valor
\[u=f(x_1,x_2,\ldots,x_n).\]
\end{frame}
% Diapositiva original 5
\begin{frame}{Conceptos de funciones de varias variables}
Los conceptos conocidos para funciones de una variable tienen su equivalente para funciones de $n$ variables.
\medskip
\begin{itemize}
\item El \textbf{dominio} contiene las entradas admitidas.
\item La \textbf{imagen} recoge los valores realmente alcanzados.
\item Las funciones pueden evaluarse, operarse y representarse.
\end{itemize}
\medskip
Para $f(x,y)=x^2+y^2$, el dominio natural es $\mathbb R^2$,
mientras que su imagen es $[0,+\infty)$.
\end{frame}
% Diapositiva original 6
\begin{frame}{Funciones de dos variables: dominio e imagen}
Particularizando a dos variables, una función $f$ de dos variables asigna a cada par ordenado de números reales $(x,y)$ perteneciente a un conjunto $D$ un número real único denotado por $u=f(x,y)$.
\medskip

El conjunto $D$ es el dominio de $f$ y su imagen es el conjunto de valores que toma $f$.
\figura{s6_0_v2.jpg}{.90}{.30}
\vspace{-2mm}
{\footnotesize En el esquema, los pares $(x,y)$ pertenecen a
$D\subseteq\mathbb R^2$; las salidas pertenecen a $\mathbb R$.
La imagen es el conjunto de las salidas obtenidas.}
\end{frame}
% Diapositiva original 7
\begin{frame}{Variables independientes y dependientes}
Las funciones de dos variables se suelen escribir de la forma $z=f(x,y)$ para explicitar el valor tomado por la función $f$ en un punto $(x,y)$.
\medskip

Las variables $x$ e $y$ son las variables independientes y $z$ es la variable dependiente.
\medskip
Por ejemplo, si $f(x,y)=x+2y$, entonces
\[
f(1,3)=7,
\]
y el punto correspondiente en la gráfica será $(1,3,7)$.
\end{frame}
% NUEVA EQUILIBRADA: Primera evaluación sencilla antes de funciones más complejas
\begin{frame}{Ejemplo básico: evaluación de una función}
Sea
\[
f(x,y)=x^2+2y.
\]
La evaluación de la función consiste en sustituir ambas variables:
\[
f(1,2)=1^2+2\cdot 2=5,\qquad
f(-1,3)=(-1)^2+2\cdot 3=7.
\]
\medskip
Al tratarse de un polinomio, $f$ está definida para todo
$(x,y)\in\mathbb R^2$: su dominio natural es $\mathbb R^2$.
\medskip
En los siguientes ejemplos aparecerán restricciones sobre las entradas.
\end{frame}

% Diapositiva original 8
\begin{frame}{Gráfica de una función de dos variables}
\small Si $f$ es una función de dos variables con dominio $D$, entonces la gráfica de $f$ es el conjunto de ternas $(x,y,z)$ de $\mathbb{R}^3$ tales que $z=f(x,y)$ y $(x,y)\in D$.
\medskip

Igual que las gráficas de funciones de una variable son curvas en el plano, las gráficas de funciones de dos variables son superficies en el espacio.
\figura{s8_2.png}{.44}{.47}
\end{frame}
% Diapositiva original 9
\begin{frame}{Ejemplo de gráfica: $z=x+y^2$}
Por ejemplo, $z=f(x,y)=x+y^2$ tiene dominio $\mathbb R^2$.
Su representación gráfica es una superficie que varía linealmente
con $x$ y cuadráticamente con $y$:
\figura{s9_2.png}{.85}{.58}
\end{frame}
% NUEVA EQUILIBRADA: Puente a los dominios originales
\begin{frame}{Determinación del dominio natural}
En los ejemplos polinómicos anteriores no hay restricciones
sobre las variables.
\medskip

En cambio, ciertas operaciones solo producen resultados reales si
sus argumentos satisfacen algunas condiciones:
\begin{itemize}
\item \textbf{Cocientes:} denominador distinto de cero.
\item \textbf{Raíces de índice par:} radicando no negativo.
\item \textbf{Logaritmos:} argumento estrictamente positivo.
\end{itemize}
\medskip
Cuando intervienen varias operaciones, debemos imponer
\textbf{todas las restricciones simultáneamente}.
\end{frame}

% Diapositiva original 10
\begin{frame}{Dominio de una función racional}
\textbf{Ejemplo.} El denominador no puede anularse.
\[z=f(x,y)=\frac{3x-5y}{y-x^2},\quad D=\{(x,y)\in\mathbb{R}^2:y-x^2\ne0\}.\]
\figura{s10_4.png}{.70}{.44}
El dominio es todo el plano $\mathbb{R}^2$ excepto los puntos de la parábola $y=x^2$.
\end{frame}
% Diapositiva original 11
\begin{frame}{Dominio de una función radical}
\textbf{Ejemplo.} El radicando debe ser no negativo.
\[z=f(x,y)=\sqrt{9-x^2-y^2},\quad D=\{(x,y)\in\mathbb{R}^2:9-x^2-y^2\ge0\}.\]
\figura{s11_5.png}{.77}{.40}
\small Como $x^2+y^2=9$ es la ecuación de una circunferencia de centro el origen y radio $3$, $D$ representa el interior y la frontera de dicha circunferencia.
\end{frame}
% Diapositiva original 12
\begin{frame}{Dominio de una función de tres variables}
El dominio de la función
\[f(x,y,z)=\ln(z-y)+xy\sin z\]
viene determinado por el logaritmo: su argumento debe cumplir
$z-y>0$. El término $xy\sin z$ no introduce restricciones.
Así pues, el dominio es
\[D=\{(x,y,z)\in\mathbb{R}^3:z>y\}.\]
Es decir, todos los puntos que están por encima del plano $z=y$.
\end{frame}
% NUEVA EQUILIBRADA: Ejemplo algo más exigente después de los dominios originales
\begin{frame}{Ejemplo adicional: restricciones combinadas}
Consideremos
\[
f(x,y)=\frac{\sqrt{4-x^2-y^2}}{y-1}.
\]
La raíz y el cociente imponen, respectivamente,
\[
4-x^2-y^2\geq0,\qquad y-1\neq0.
\]
Por tanto, el dominio natural es
\[
\boxed{D=\{(x,y)\in\mathbb R^2:
x^2+y^2\leq4,\ y\neq1\}.}
\]
\medskip
Es el disco cerrado de radio $2$, excluyendo los puntos
que pertenecen a la recta $y=1$.
\end{frame}

% Diapositiva original 13
\begin{frame}{Operaciones con funciones de varias variables}
Las funciones de varias variables pueden operarse de igual forma que las de una variable, ya sea con suma, producto, cociente o composición de funciones.
\medskip
Por ejemplo, si $f(x,y)=x+y$ y $g(x,y)=x^2+y^2$,
\[
(f+g)(x,y)=x+y+x^2+y^2,
\]
\[
\left(\frac fg\right)(x,y)=\frac{x+y}{x^2+y^2}.
\]
\small
La suma tiene dominio natural $\mathbb R^2$. En el cociente
es preciso excluir el punto $(0,0)$, donde $g=0$.
\end{frame}
% NUEVA EQUILIBRADA: Contexto para el siguiente apartado
\begin{frame}{De las superficies a las curvas de nivel}
La gráfica de una función de dos variables es, en general,
una superficie del espacio tridimensional.
\medskip

En cartografía y visualización científica resulta útil describir
esa superficie mediante información representada en un plano.
\medskip

Para hacerlo, fijamos un valor $k$ y buscamos todos los puntos
del dominio en los que
\[
f(x,y)=k.
\]
\medskip
Estos conjuntos se denominan \textbf{curvas de nivel} cuando
su geometría es la de una curva.
\end{frame}

% Diapositiva original 14
\begin{frame}{Curvas de nivel}
Las curvas de nivel de una función $f$ de dos variables son las curvas con ecuaciones $f(x,y)=k$, donde $k$ es una constante que pertenece a la imagen de $f$.
\begin{columns}[c]
\begin{column}{.48\textwidth}\figura{s14_3.png}{.85}{.56}\end{column}
\begin{column}{.48\textwidth}\figura{s14_5.png}{.98}{.56}\end{column}
\end{columns}
\end{frame}
% NUEVA EQUILIBRADA: Ejemplo muy sencillo tras la definición formal
\begin{frame}{Ejemplo introductorio: curvas de nivel rectas}
Sea la función $f(x,y)=x+y$.
Para obtener su curva de nivel correspondiente al valor $k$:
\[
f(x,y)=k\quad\Longleftrightarrow\quad y=k-x.
\]
\medskip
Todas las curvas de nivel son \textbf{rectas paralelas}:
\[
\begin{array}{c|c}
k&\text{ecuación del nivel}\\\hline
-1&y=-1-x\\
0&y=-x\\
1&y=1-x
\end{array}
\]
\medskip
A continuación conservamos el ejemplo original, cuyas curvas
de nivel son circunferencias.
\end{frame}

% Diapositiva original 15
\begin{frame}{Curvas de nivel: ejemplo de una semiesfera}
\textbf{Ejemplo}
\vfill
Obtener las curvas de nivel de
\[z=\sqrt{100-x^2-y^2}.\]
\medskip
El dominio natural es el disco $x^2+y^2\leq100$.
Se trata de estudiar qué posiciones corresponden a una
\textbf{altura constante} $z=c$.
\vfill
\end{frame}
% Diapositiva original 16
\begin{frame}{Resolución: curvas de nivel de la semiesfera}
\small\textbf{Ejemplo}\quad $z=\sqrt{100-x^2-y^2}$.
\medskip

Sustituimos la variable $z$ por una constante $z=c$:
\[c=\sqrt{100-x^2-y^2}.\]
Elevando al cuadrado,
\[c^2=100-x^2-y^2.\]
Reorganizando,
\[x^2+y^2=100-c^2.\]
Se obtienen circunferencias de centro el origen para $0\le c<10$. Para $c=10$ se obtiene el punto $(0,0)$. Para $c<0$ o $c>10$ no hay curvas de nivel.
\end{frame}
% Diapositiva original 17
\begin{frame}{Representación de los niveles obtenidos}
\textbf{Interpretación gráfica}\par
Representar los niveles de la función de dos variables
\begin{columns}[c]
\begin{column}{.43\textwidth}
\[z=\sqrt{100-x^2-y^2}\]
\[x^2+y^2=100-c^2\]
\end{column}
\begin{column}{.55\textwidth}\figura{s17_3.png}{1}{.58}\end{column}
\end{columns}
\end{frame}
% Diapositiva original 18
\begin{frame}{Superficie y curvas de nivel}
\figura{s18_0.jpg}{.78}{.83}
\end{frame}
% NUEVA EQUILIBRADA: Explicación de la primera imagen clásica
\begin{frame}{Interpretación: cortes y proyecciones}
La figura anterior relaciona una superficie con sus curvas de nivel.
\medskip
\begin{enumerate}
\item Fijamos una altura $z=k$ y consideramos el plano horizontal correspondiente.
\medskip
\item La intersección de ese plano con la superficie identifica los puntos que tienen el mismo valor $k$.
\medskip
\item Al proyectar estos puntos sobre el plano $XY$ obtenemos la curva de nivel.
\end{enumerate}
\medskip
Variando $k$ se construye una representación plana del relieve
de la superficie.
\end{frame}

% Diapositiva original 19
\begin{frame}{Curvas de nivel: una función no radial}
\figura{s19_1.jpg}{1}{.72}
\begin{center}\scriptsize Dos vistas de $f(x,y)=-xy\,e^{-x^2-y^2}$ y sus curvas de nivel.\end{center}
\end{frame}
% NUEVA EQUILIBRADA: Explicación de la segunda figura clásica
\begin{frame}{Interpretación: regiones positivas y negativas}
La figura anterior representa
\[
f(x,y)=-xy\,e^{-x^2-y^2}.
\]
\medskip
Como $e^{-x^2-y^2}>0$ para todo $(x,y)$, el signo depende de $-xy$:
\[
\begin{array}{ll}
f(x,y)>0,&xy<0,\\
f(x,y)<0,&xy>0,\\
f(x,y)=0,&xy=0.
\end{array}
\]
\medskip
Esto explica la aparición de \textbf{dos regiones positivas}
y \textbf{dos negativas}. Las curvas de nivel no tienen por qué
ser circunferencias, ni siquiera un único lazo cerrado.
\end{frame}

% Diapositiva original 20
\begin{frame}{Curvas de nivel: una función racional}
\figura{s20_1.jpg}{1}{.76}
\end{frame}
% NUEVA EQUILIBRADA: Demostración de las curvas de la tercera figura clásica
\begin{frame}{Cálculo de los niveles de la función racional}
Para la función de la figura anterior,
\[
f(x,y)=\frac{-3y}{x^2+y^2+1},
\]
fijamos un nivel $k\ne0$ y despejamos:
\[
k(x^2+y^2+1)=-3y,
\]
\[
\boxed{x^2+\left(y+\frac{3}{2k}\right)^2
=\frac{9}{4k^2}-1.}
\]
\small
Para $0<|k|<\frac32$ se obtienen circunferencias;
para $|k|=\frac32$, un punto; para $|k|>\frac32$, ninguno.
Si $k=0$, la curva de nivel es la recta $y=0$.
\end{frame}

% Diapositiva original 21
\begin{frame}{Aplicación: mapa topográfico}
\figura{s21_0.jpg}{.73}{.86}
\end{frame}
% NUEVA EQUILIBRADA: Interpretación de la cuarta imagen clásica
\begin{frame}{Interpretación: lectura de un mapa topográfico}
El mapa anterior representa la altura mediante curvas de nivel:
cada curva reúne posiciones con \textbf{igual cota}.
\medskip
\begin{itemize}
\item Las curvas cerradas alrededor de una cumbre describen
distintas alturas del terreno.
\medskip
\item A igualdad de intervalo de cotas, cuanto \textbf{más próximas}
estén las curvas, mayor será la variación de altura por unidad
de distancia horizontal.
\medskip
\item El mismo principio se utiliza para visualizar datos
espaciales y campos escalares.
\end{itemize}
\end{frame}

% NUEVA EQUILIBRADA: Una conexión informática breve, no un hilo conductor
\begin{frame}{Aplicación informática: funciones de coste}
En optimización y aprendizaje automático es habitual estudiar
funciones que dependen de varios parámetros.
\medskip
Por ejemplo,
\[
L(w_1,w_2)=(w_1-2)^2+2(w_2+1)^2.
\]
Sus curvas de nivel verifican
\[
(w_1-2)^2+2(w_2+1)^2=k.
\]
\medskip
Si $k>0$, se obtienen \textbf{elipses concéntricas};
el mínimo, $L=0$, se alcanza en $(2,-1)$.
Estas curvas ayudan a visualizar la forma de una función objetivo.
\end{frame}

% NUEVA EQUILIBRADA: Práctica complementaria de primer nivel
\begin{frame}{Ejercicios complementarios: nivel básico}
\begin{enumerate}
\item Evaluar $f(x,y)=x^2+y$ en $(1,2)$ y $(-2,3)$.
\medskip
\item Encontrar el dominio natural de
\[
g(x,y)=\sqrt{4-x^2-y^2}.
\]
\item Describir las curvas de nivel para $k=0,1,2$ de
\[
h(x,y)=x+y.
\]
\end{enumerate}
\medskip
\small
Objetivo: aplicar directamente las definiciones y procedimientos.
\end{frame}

% NUEVA EQUILIBRADA: Práctica complementaria de segundo nivel
\begin{frame}{Ejercicios complementarios: nivel intermedio}
\begin{enumerate}
\item Determinar y representar el dominio de
\[
f(x,y)=\ln(y-x^2).
\]
\item Encontrar el dominio de
\[
g(x,y)=\frac{\sqrt{4-x^2-y^2}}{y-1}.
\]
\item Describir los niveles $k=-1,0,1$ de
\[
h(x,y)=x^2-y^2.
\]
\end{enumerate}
\medskip
\small
Los ejercicios originales que siguen incluyen problemas de mayor integración.
\end{frame}

% Diapositiva original 22
\begin{frame}{Ejercicios originales: consolidación y ampliación}
\small\textbf{Ejercicios}
\medskip

Encontrar el dominio de la función
\[f(x,y,z)=\arcsin(xy)\,e^{3z}+e^{(y+z)\ln(x+z)}.\]
Describir los conjuntos de nivel para los valores de $k$ indicados:
\begin{enumerate}[a)]
\item $f(x,y)=x^2+y^2$,\quad $k=0,1,2,3$.
\item $f(x,y,z)=x^2+y^2$,\quad $k=0,1,2$.
\item $f(x,y)=x^2-y^2$,\quad $k=-2,-1,0,1,2$.
\end{enumerate}
En a) y c) se trata de curvas de nivel; en b), de superficies de nivel.
\end{frame}
\end{document}
